\documentclass[11pt]{article}
\usepackage[a4paper,margin=1in]{geometry}
\usepackage{amsmath,amssymb,amsthm,mathtools}
\usepackage{enumitem}
\usepackage{microtype}
\usepackage[colorlinks=true,linkcolor=blue,citecolor=blue,urlcolor=blue]{hyperref}
\usepackage[nameinlink,capitalise,noabbrev]{cleveref}

\newtheorem{theorem}{Theorem}[section]
\newtheorem{lemma}[theorem]{Lemma}
\newtheorem{proposition}[theorem]{Proposition}
\newtheorem{corollary}[theorem]{Corollary}
\newtheorem{definition}[theorem]{Definition}
\newtheorem{problem}[theorem]{Open Problem}
\newtheorem{remark}[theorem]{Remark}

\newcommand{\cP}{\mathcal P}
\newcommand{\cF}{\mathcal F}
\newcommand{\IT}{\mathrm{IT}}
\newcommand{\supp}{\operatorname{supp}}
\newcommand{\defect}{\operatorname{def}}
\newcommand{\one}{\mathbf 1}

\title{A Rank-Two Repair Program for Independent Transversals\\
       in Partitioned Three-Uniform Hypergraphs}
\author{}
\date{6 August 2026}

\begin{document}
\maketitle

\section{Setting and current status}

Let \(H=(V,E)\) be a three-uniform hypergraph.  Its vertex set is partitioned
into blocks
\[
  \cP=\{B_1,\dots,B_m\},
  \qquad |B_i|=b.
\]
Every edge is assumed to meet each block in at most one vertex.  A
\emph{transversal} is a set containing exactly one vertex from every block, and
an \emph{independent transversal} is a transversal containing no edge of
\(H\).  Write \(d_H(v)\) for the degree of \(v\) and
\(\Delta(H)=\max_v d_H(v)\).

The target assertion is
\begin{equation}\label{eq:target}
  \Delta(H)<\left(\frac14-o(1)\right)b^2
  \quad\Longrightarrow\quad
  H\text{ has an independent transversal}.
\end{equation}
The assertion remains open.  The purpose of this paper is to isolate a route
whose local counting naturally yields the constant \(1/4\), to prove the
available local and conditional structural statements, and to state the one
remaining global decomposition precisely.

The route has four components:
\[
\boxed{
\begin{aligned}
&\text{finite actual-history organization}\longrightarrow
  \text{rank-two repair},\\
&\text{rank-two repair}\longrightarrow
  \text{positive disjoint-blocker fork density},\\
&\text{fork density}\longrightarrow
  \text{inverse-fiber codimension or a persistent anchor},\\
&\text{persistent anchor}\longrightarrow d_H(p)\ge b^2.
\end{aligned}}
\]
The unresolved step is the third arrow in its global, history-preserving form.

\section{Blockers and canonical maximal-matching repair}

A \emph{partial transversal} is a set meeting every block in at most one
vertex.  Let \(S\) be an independent partial transversal and let \(x\) belong
to a block missed by \(S\).  Since \(S\) is independent, every edge contained
in \(S\cup\{x\}\) contains \(x\).

\begin{definition}[Blocker-link graph]
The blocker-link graph at \((S,x)\) is
\[
  G_x(S)=\bigl\{\{u,v\}\subseteq S:\{x,u,v\}\in E(H)\bigr\}.
\]
Fix an order on all vertex pairs and let \(M_x(S)\) be the greedy maximal
matching of \(G_x(S)\) in that order.  Put
\[
  r_x(S)=|M_x(S)|.
\]
\end{definition}

\begin{lemma}[Canonical matching repair]\label{lem:matching-repair}
The set
\[
  R_x(S)=(S\cup\{x\})\setminus V(M_x(S))
\]
is an independent partial transversal.  Moreover,
\[
  |R_x(S)|-|S|=1-2r_x(S).
\]
\end{lemma}

\begin{proof}
If a blocker pair \(f\in G_x(S)\) were disjoint from \(V(M_x(S))\), then
\(f\) could be added to the matching, contradicting maximality.  Thus
\(V(M_x(S))\) meets every blocker pair, and deleting it destroys every edge
created by adding \(x\).  The block condition is preserved because only old
vertices are deleted and \(x\) lies in a missing block.  The size identity is
immediate.
\end{proof}

\begin{definition}[Disjoint-blocker fork]
A repair occurrence is a disjoint-blocker fork if \(r_x(S)\ge2\).  Choosing
the first two matching edges gives actual hyperedges
\[
  e_1=\{x,u_1,v_1\},\qquad e_2=\{x,u_2,v_2\}
\]
with
\[
  \{u_1,v_1\}\cap\{u_2,v_2\}=\varnothing.
\]
The quantity \((r_x(S)-1)_+\) is the fork excess of the occurrence.
\end{definition}

A blocker family may contain many edges while its matching number is one.  In
that case two deleted old vertices hit the entire family.  Thus the relevant
obstruction is not ``more than one blocker'' but matching number at least two.

\section{Record growth and the constant one quarter}

We use a standard reconstruction hypothesis.  A deterministic repair protocol
is called \emph{faithfully reconstructible} if its finite execution record,
final partial transversal, and fixed initial data determine the sequence of
attempted vertices.  Actual edge labels and all choices made by the canonical
matching rule are retained in the record.

\subsection{Fork-free records}

Assume first that \(r_x(S)\in\{0,1\}\) at every step.  A successful insertion
has height change \(+1\) and no edge label.  A rank-two repair has height change
\(-1\) and records one edge containing the attempted vertex, with at most
\(\Delta(H)\) choices.  The node polynomial is therefore
\[
  \Phi_0(z)=1+\Delta(H)z^2.
\]

\begin{proposition}[Fork-free growth]\label{prop:fork-free}
For a faithfully reconstructible fork-free protocol, the exponential growth
rate of records is at most
\[
  \inf_{z>0}\frac{1+\Delta(H)z^2}{z}=2\sqrt{\Delta(H)}.
\]
Consequently, if every word in an alphabet of size \(b\) generated a
nonterminating fork-free record, then \(b\le2\sqrt{\Delta(H)}\).
\end{proposition}

\begin{proof}
The weighted tree bound gives the coefficient growth controlled by
\(\inf_{z>0}\Phi_0(z)/z\).  The minimum of \(z^{-1}+\Delta z\) occurs at
\(z=\Delta^{-1/2}\).  Faithful reconstruction injects the choice words into
records up to subexponential initial and terminal data, yielding the stated
inequality.
\end{proof}

\begin{corollary}\label{cor:fork-free-quarter}
Within the faithfully reconstructible fork-free class,
\[
  \Delta(H)<\frac{b^2}{4}
\]
forces termination at an independent transversal.
\end{corollary}

This is the first important reason to use maximal-matching repair: the target
constant appears before any global stability or terminal accounting.

\subsection{Fork density below the threshold}

Introduce a weight \(y^{(r-1)_+}\), where \(0<y<1\).  Bounding the number of
possible ordered matching-edge labels by \(\Delta^r\) gives
\[
  \Phi_y(z)
  =1+\sum_{r\ge1}y^{r-1}\Delta^r z^{2r}
  =1+\frac{\Delta z^2}{1-y\Delta z^2}.
\]
At \(z=\Delta^{-1/2}\),
\begin{equation}\label{eq:weighted-growth}
  \frac{\Phi_y(z)}{z}
  =\sqrt\Delta\,\frac{2-y}{1-y}.
\end{equation}

Fix \(0<\varepsilon<1/4\), set \(y=\varepsilon\), and define
\[
  A_\varepsilon
  =\frac{\sqrt{1-4\varepsilon}}2
    \frac{2-\varepsilon}{1-\varepsilon}.
\]
A direct calculation gives
\[
  4(1-\varepsilon)^2
  -(1-4\varepsilon)(2-\varepsilon)^2
  =\varepsilon(4\varepsilon^2-13\varepsilon+12)>0,
\]
so \(A_\varepsilon<1\).  Put
\[
  \alpha_\varepsilon
  =\frac{-\log A_\varepsilon}{2\log(1/\varepsilon)}>0.
\]

\begin{proposition}[Positive fork density]\label{prop:fork-density}
Suppose
\[
  \Delta(H)\le\left(\frac14-\varepsilon\right)b^2.
\]
In a faithfully reconstructible nonterminating protocol, records of length
\(T\) with total fork excess at most \(\alpha_\varepsilon T\) have exponential
growth strictly smaller than \(b^T\).  Hence, if all \(b\)-ary choice words
avoid an independent transversal, an exponentially dominant set of long
records has positive linear fork excess.
\end{proposition}

\begin{proof}
If the excess is at most \(\alpha T\), then
\(y^{\mathrm{excess}}\ge y^{\alpha T}\).  Dividing the weighted record count by
this factor and using \cref{eq:weighted-growth} bounds the growth by
\[
  y^{-\alpha}\sqrt\Delta\frac{2-y}{1-y}.
\]
With the choices above and
\(\sqrt\Delta\le\tfrac12\sqrt{1-4\varepsilon}\,b\), the resulting base is
\(\sqrt{A_\varepsilon}\,b<b\).
\end{proof}

\section{Finite actual-history linear programming}

The entropy argument produces fork occurrences.  A separate issue is the
organization of long histories.  Let \(D=(\Omega,A)\) be a finite directed
graph of actual history states after all already accepted exits have been
deleted.

\begin{proposition}[Potential--core equivalence]\label{prop:potential-core}
The following are equivalent:
\begin{enumerate}[label=(\roman*)]
\item there is a function \(\varphi:\Omega\to\mathbb R\) that strictly
decreases on every arc of \(D\);
\item \(D\) is acyclic;
\item \(D\) has no nonzero nonnegative circulation.
\end{enumerate}
If these conditions fail, some bottom strongly connected component supports a
nonzero circulation.
\end{proposition}

\begin{proof}
A strict potential excludes directed cycles.  Every finite acyclic digraph has
a topological ordering, which is a strict potential after reversing sign.  A
directed cycle supports a nonzero circulation, while a nonzero nonnegative
circulation contains a directed cycle in its positive support.
\end{proof}

This proposition is an exhaustion tool.  It does not imply that a recurrent
class has bounded inverse multiplicity, that it is closed under every legal
release omitted by a deterministic policy, or that it produces a negative
term.  The state space must retain all history needed to determine successor
classification.

\section{Fork loads and inverse fibers}

Fix a pivot \(x\).  Let \(w_x(e,f)\ge0\) be a weight on unordered pairs of
actual edges through \(x\) whose non-pivot endpoint pairs are disjoint.  Define
\[
  F_x=\sum_{\{e,f\}}w_x(e,f),
  \qquad
  \ell_x(e)=\sum_f w_x(e,f),
  \qquad
  \ell_x^*=\max_{e\ni x}\ell_x(e).
\]

\begin{lemma}[Fork load double count]\label{lem:fork-load}
\[
  2F_x=\sum_{e\ni x}\ell_x(e)
  \le d_H(x)\ell_x^*.
\]
\end{lemma}

Thus a bound
\[
  \ell_x^*\le\frac{8+o_b(1)}{b^2}F_x
\]
would give
\[
  d_H(x)\ge\left(\frac14-o_b(1)\right)b^2.
\]
The issue is therefore not the existence of many forks but their possible
reuse of a small set of actual edges.

\subsection{Two-coordinate replacement boxes}

Assume that \(H\) is edge-minimal subject to having no independent
transversal.  For every edge \(e\), there is a full transversal \(W_e\) with
\[
  E(H[W_e])=\{e\}.
\]
Indeed, an independent transversal of \(H-e\) must contain \(e\).

Fix \(e=\{x,a,b\}\), where \(a\in A\) and \(b\in B\).  For
\(u\in A\), \(v\in B\), set
\[
  W_e(u,v)=W_e-a-b+u+v.
\]
Since \(H\) has no independent transversal, choose an edge
\(g_e(u,v)\subseteq W_e(u,v)\) by a fixed rule.

\begin{proposition}[Replacement-box codimension]\label{prop:box}
Every \(g_e(u,v)\) contains \(u\) or \(v\).  For a fixed actual edge \(g\),
\[
\left|\{(u,v):g_e(u,v)=g,\ \{u,v\}\subseteq g\}\right|\le1,
\]
and
\[
\left|\{(u,v):g_e(u,v)=g,\ |g\cap\{u,v\}|=1\}\right|\le b.
\]
The same estimates hold for weighted parent occurrences when the parent
history is retained as part of the sample point.
\end{proposition}

\begin{proof}
If \(g_e(u,v)\) contains neither replacement coordinate, then it is contained
in \(W_e\setminus\{a,b\}\), contradicting
\(E(H[W_e])=\{e\}\).  If \(g\) contains both replacements, its vertices in
blocks \(A,B\) determine \((u,v)\).  If it contains exactly one, that coordinate
is determined and the other has at most \(b\) choices.  Retaining the parent
occurrence makes the weighted statement a disjoint sum of these pointwise
bounds.
\end{proof}

The first class has the desired \(b^{-2}\) inverse multiplicity.  The second
has only \(b^{-1}\).  An output in that class can have the form
\[
  g_e(u,v)=\{u,p,q\},
\]
where the pair \(\{p,q\}\) is fixed along one row but may change from row to
row.  Such row-dependent coverings can occur with maximum degree of order
\(b\); therefore local box counting alone cannot produce a persistent anchor.

\section{The persistent-anchor endgame}

Let \(p\) be a vertex.  We formulate the exact closure hypothesis needed for
target-following.

\begin{definition}[Future-complete fixed-pivot class]
A class \(\mathcal C_p\) consists of independent one-hole transversals and is
future-complete for \(p\) if, whenever a state in \(\mathcal C_p\) attempts a
vertex in its missing block, the following hold unless an already accepted exit
occurs:
\begin{enumerate}[label=(\alph*)]
\item there is a unique blocker;
\item that blocker contains \(p\);
\item every legal release successor that preserves the fixed-pivot procedure
belongs to \(\mathcal C_p\).
\end{enumerate}
\end{definition}

\begin{theorem}[Fixed-pivot target covering]\label{thm:anchor}
If a nonempty class \(\mathcal C_p\) is future-complete for \(p\) and has no
accepted exit, then
\[
  d_H(p)\ge b^2.
\]
\end{theorem}

\begin{proof}
Fix a target \(Q\) choosing one vertex from every block other than the block of
\(p\).  Starting from a state in \(\mathcal C_p\), attempt the target vertex in
the missing block.  The unique blocker has the form
\(\{p,q_M,z\}\), where \(q_M\) is the attempted target vertex and \(z\) is an
old selected vertex.  Release \(z\).

Let \(\Psi\) be the number of non-pivot blocks whose selected vertex agrees
with \(Q\).  If \(z\) is not the target vertex of its block, then the operation
installs \(q_M\) and deletes a non-target vertex, so \(\Psi\) increases.  If
\(z\) is the target vertex of its block, then \(Q\) contains the link edge
\(\{q_M,z\}\in E(L_H(p))\).  Since \(\Psi\) is bounded, every target contains a
link edge of \(p\).

There are \(b^{m-1}\) targets outside the block of \(p\).  Each link edge fixes
two block coordinates and belongs to exactly \(b^{m-3}\) targets.  Hence
\[
  b^{m-1}\le d_H(p)b^{m-3},
\]
which proves the claim.
\end{proof}

The hypothesis of \cref{thm:anchor} is genuinely global.  A heavy vertex or
pair in one inverse fiber need not survive the next release, need not occur in
the next blocker, and need not be shared by different branches.

\section{The active global decomposition problem}

We can now state the missing theorem without referring to an artificial
terminal category.

\begin{problem}[Global inverse-fiber decomposition]\label{prob:global-fiber}
For the weighted actual fork occurrences produced by
\cref{prop:fork-density}, construct a decomposition, preserving parent history,
owner, root and actual edge identities, into the following four disjoint parts:
\begin{enumerate}[label=(\roman*)]
\item a two-coordinate diffuse part satisfying
\[
  \ell_x^*\le\frac{8+o_b(1)}{b^2}F_x;
\]
\item a part reaching an independently verified exit or strict progress term;
\item a part contained in a no-IT subsystem induced by a proper set of complete
blocks;
\item a positive-mass future-complete class with one fixed persistent anchor.
\end{enumerate}
The decomposition and its constants must be uniform in the number of blocks and
in the size of the finite history graph.
\end{problem}

The two endpoints are already available.  Part~(i) combines with
\cref{lem:fork-load}; part~(iv) combines with \cref{thm:anchor}.  The precise
unresolved behavior is one-coordinate concentration whose heavy pair migrates
when an external coordinate is first lost.  A proof should therefore be
organized by the first anchor-loss time and an inclusion-minimal candidate
anchor.

\section{Why the finite stopping route is not a closure theorem}

A local three-coordinate contraction can be iterated in fully unfolded history
space, but four distinctions are essential.

First, an independent completion on the currently exposed proper block set is
not a global independent transversal.  Its terminal mass cannot be deleted
from a stopping identity using only the global no-IT assumption.

Second, whether a successor is a fresh state or a return depends on the visited
history.  A label consisting only of the current physical configuration,
blocker and seen resources is not a transition congruence.

Third, a return along one deterministic release policy gives only a policy
cycle.  It does not imply closure under every legal release, which is required
by an all-release recurrent-core theorem.

Fourth, a tail bound whose depth depends on
\(|E(H)|+|V(H)|\) is a fixed-instance statement.  It cannot be inserted as a
uniform asymptotic error without a separate quantitative argument.

These observations do not invalidate finite history organization; they limit
its role to the first stage of the program.

\section{Assessment of competing routes}

The rank-two route is distinguished by three features.

\begin{enumerate}[label=(\arabic*)]
\item It produces the constant \(1/4\) directly from the repair polynomial
\(1+\Delta z^2\).
\item It treats high inverse multiplicity as a possible source of rigidity,
rather than assuming bounded multiplicity.
\item Its zero-loss branch has a concrete endgame: a fixed pivot whose link
covers every target, forcing degree at least \(b^2\).
\end{enumerate}

Finite-state LP, root normalization and critical-stability identities remain
valuable for exhaustion and bookkeeping.  They do not by themselves convert
residual circulation into a degree lower bound.  Global signed re-entry, Hall
transport, waiting-time coding, first-owner compression, stationary matching,
and private-target external-coordinate concentration all fail at that
conversion unless supplemented by a new actual-object theorem.  The active
problem \cref{prob:global-fiber} is therefore the most concentrated remaining
bridge among the explored approaches.

\section{Conditional closure}

\begin{theorem}[Conditional one-quarter conclusion]
Assume that the global inverse-fiber decomposition in
\cref{prob:global-fiber} holds with the stated uniform constants, and that the
proper-block branch is excluded by block minimality.  Then, for every fixed
\(\varepsilon>0\) and all sufficiently large \(b\),
\[
  \Delta(H)\le\left(\frac14-\varepsilon\right)b^2
\]
forces an independent transversal.
\end{theorem}

\begin{proof}
If no independent transversal exists, \cref{prop:fork-density} yields positive
fork mass.  In the diffuse branch, \cref{lem:fork-load} and the required
\(b^{-2}\) inverse-multiplicity estimate give a vertex of degree
\((1/4-o_b(1))b^2\), contradicting the assumed gap for large \(b\).  The exit
branch is absorbed by its independent progress statement, the proper-block
branch contradicts minimality, and the anchor branch contradicts
\cref{thm:anchor}.  Hence the no-IT assumption is impossible.
\end{proof}

\end{document}
